\section{Limits of recoverability}
\label{sec:scope}
%==========================================================================

\emph{This section is a scope limit on everything above, stated here rather than in a limitations paragraph because it bites hardest in the population the question matters
most for.}

In \eqref{eq:local} the alternative channel does not saturate, so as residual capacity
falls the marker does not rise without bound. It approaches a ceiling set by that channel
alone,
\begin{equation}
m \;\longrightarrow\; \frac{L_0 + L}{k_e} \qquad \text{as } \theta \to 0,
\label{eq:ceilinglimit}
\end{equation}
which at the illustrative values is $6.45$ times the managed level. The elasticity of the
marker to capacity dies with it.

\begin{center}
\footnotesize
\begin{tabular}{@{}rrr@{}}
\toprule
$\theta$, as \% of nominal & marker, $\times$ managed & $\lvert\partial \ln m / \partial \ln\theta\rvert$ \\
\midrule
$100\%$ & $1.00$ & $0.915$ \\
$50\%$  & $1.82$ & $0.807$ \\
$20\%$  & $3.41$ & $0.537$ \\
$10\%$  & $4.57$ & $0.321$ \\
$5\%$   & $5.40$ & $0.173$ \\
$2\%$   & $6.00$ & $0.072$ \\
$1\%$   & $6.22$ & $0.036$ \\
\bottomrule
\end{tabular}
\end{center}

Stated in the units a clinic reads: once a patient's managed marker sits above
about three to four times the reference level, residual capacity is no longer recoverable
from it. The elasticity falls below $0.5$ at $3.6$ times managed, below $0.25$ at $5.0$,
and below $0.1$ at $5.8$. At the bottom of that table two patients differing tenfold in
capacity, $10\%$ of nominal against $1\%$, are separated by less than a third in
marker level, which no assay resolves against the day-to-day variability of the quantity.

\subsection{What the estimate's own precision does across that range}
\label{sec:scopeprec}

The table above is about the \emph{marker}. The same question can be put to the
\emph{estimate}, and until now it had not been: every precision this paper reports is
computed at $\theta = 3000$, and on the table above that is the $100\%$ row.

That is worth stating plainly, because the repository had left it ambiguous.
\texttt{src/common/models.py} annotates $\theta = 3000$ as $V_{\max}$, the \emph{maximal}
enzymatic velocity, and annotates the parent model's $\theta = 300$ as ``$\sim 10\%$ of
normal''. Together those put normal at $3000$ and the local block at $100\%$ of it, so the
local block is a normal-capacity illustration and Table~\ref{tab:params} now labels it as
one. Residual PCC activity in propionic acidemia is under $5\%$ of
control~\cite{pcc_residual_activity}, so the patient sits at $\theta \le 150$, which the
table now carries as its own row. Nothing about the model changes; what changes is that the
illustration is no longer available to be read as a description of the patient.

Holding the recorded load fixed and letting the marker rise, which is the convention the
table above already uses, and re-solving the operating level, $\sigma$, the load-bias prior
and the excitation at each capacity:

\begin{center}
\footnotesize
\begin{tabular}{@{}rrrrr@{}}
\toprule
$\theta$, as \% of nominal & marker, $\times$ managed &
$\lvert\partial \ln m / \partial \ln\theta\rvert$ &
$\operatorname{sd}(\ln\hat\theta)$ & monthly visits to $0.20$ \\
\midrule
$100\%$ & $1.00$ & $0.915$ & $\mathbf{0.1699}$ & $11$ \\
$20\%$  & $3.41$ & $0.537$ & $0.8107$ & --- \\
$10\%$  & $4.57$ & $0.321$ & $1.5830$ & --- \\
$5\%$   & $5.40$ & $0.173$ & $3.0673$ & --- \\
$2\%$   & $6.00$ & $0.072$ & $7.4572$ & --- \\
\bottomrule
\end{tabular}
\end{center}

\noindent
Medians over $24$ simulated intake records at the specification's own $10\%$ day-to-day
variability; the seed ranges are in \texttt{audit/n4\_theta\_scope.py}. A dash means the
target is not reached at any number of monthly visits up to $600$, which is fifty years of
follow-up.

$\operatorname{sd}(\ln\hat\theta)$ is not monotone in $\theta$, so the gate is crossed twice.
It holds for $\theta$ between $2472$ and $18814$, which is $82\%$ to $627\%$ of nominal, a
band of $7.6$ fold, with the best precision $0.123$ at about twice nominal. Below the band
the marker has stopped responding to capacity, which is the statement above. Above it the
$\theta$ column saturates at $L$ while the assay term keeps growing, so information per draw
stops improving. \textbf{The disease sits a factor of $16$ below the lower edge of the band
the design works in.} Sampling every two days for two years instead of monthly, fifteen times
the draws, moves that edge only from $82\%$ to $75\%$ of nominal. That $9\%$ is the
quantitative form of the sentence below: it is not a statement about sampling effort.

One consequence follows for the dose and it is independent of the precision. The
admissibility bound $\theta + \Delta_{\mathrm{peak}} \le$ normal is satisfied at the
illustrative point only because the check reads $\theta$ as the \emph{residual} and infers
normal as $\theta/0.05$; written out, its verdict is
$\operatorname{frac} \times (1 + \Delta_{\mathrm{peak}}/\theta)$ with $\theta$ cancelled,
which cannot report a breach below a residual activity of $32\%$ and so cannot bind on this
disease at all. Read instead with $\theta = 3000$ as normal and $D$ held at its booked value
rather than scaled with capacity, which is what a dose in mg/kg is, the bound is breached at
$10\%$ of nominal and below: at $5\%$ the booked increment would drive the enzyme to $2.18$
times normal, and the largest admissible increment there is $D = 2672$ rather than $6000$.
That does not rescue the precision, since at $D = 2672$ and $\theta = 150$ the monthly design
returns $2.96$, but it is a second and independent reason to think the increment is set too
high, alongside the trial's own reported falls in the $D$ row of Table~\ref{tab:params}.

This is the severely affected patient. The population in which the dosing question
is most urgent is the population in which the record cannot answer it. This is not a statement about sampling effort: no schedule, no excursion and no precision on the dose
recovers a quantity the observable has stopped responding to. The therapy channel of
Section~\ref{sec:excitation} does not rescue it either: it supplies excitation in
$\theta$ against $K_m$, not sensitivity of $m$ to $\theta$.

One consequence is favourable and should not be lost in the statement of a
limit. What dies here is the estimate of \emph{capacity}. A prediction read directly off
the marker does not die, because the marker is what is measured. So the failure is
asymmetric: it falls on the mechanistic quantity and not on the observable, which is an
argument for keeping claims close to the observable in exactly this population.

\vspace{0.5em}
\hrule
\vspace{1em}

%==========================================================================
